% abstract.tex
We reproduce and extend the pLDDT-fragmentation analysis (Q0/Q4) of
Cazals~\&~Sarti~\cite{cazals2025alphafold} at the scale of the human
(\textit{H.~sapiens}) proteome using AlphaFold Database v4 predictions.
The analysis is grounded in a path-graph filtration of residues ordered by
decreasing predicted local-distance difference test (pLDDT) confidence
score, from which a per-protein persistence diagram is derived using
Union-Find with the Elder Rule.
Five statistics are computed from each diagram: the fraction of
positive-persistence pairs~$f^+_{cp}$, mean persistence~$\bar{p}$,
normalised persistence entropy~$H_p$, the peak connected-component
count~$N_{cc}^{max}$, and the count of persistent local maxima~$\mathrm{PLM}(t_p)$.

From all 23{,}391 fragments we recover a Pearson correlation of
$r = 0.864$ between $f^+_{cp}$ and~$H_p$ (paper target:
${\approx}0.97$); the residual gap is traced
to float-precision pLDDT values in v4, which produce substantially more
distinct filtration levels per protein than the effectively integer-valued
pLDDT used in the paper.
At threshold $t_p = 0.025$ and the paper's three-way filter
($n \ge 200$, $H_p \ge 0.25$, $\mathrm{PLM} \ge 3$),
we identify 183 multi-domain fragmentation candidates
(paper target: ${\approx}86$), again explained by the same precision shift.

As a novel extension we compute $\mathrm{C}_\alpha$-packing arity
signatures~\cite{cazals2025alphafold} for the entire proteome and
cross-correlate them with the fragmentation signal,
finding a 5.92-fold enrichment of fragmented proteins in the arity bin
($a_{25}=7$, $a_{75}=15$; $p = 1.61 \times 10^{-3}$), suggesting
that multi-domain proteins with compact core packing are the primary
fragmentation targets.
All code is written in Python~3.11 and released with the submission.
